%! Choosing a Sample: Four Sampling Techniques — Unit 1, Lesson 1.3 — Guided Notes & Practice
% THE in-class document, in the MAIN IDEAS / NOTES shape (LESSON_SHAPE.md §1, ported from AP
% Statistics 2026-09-12; 1.3 converted 2026-09-13 from the pilot's boxed sections): the
% vocabulary box, the hook, then ONE two-column guidednotes table. Each row is one idea — a
% label on the left; on the right one or two COMPLETE printed sentences, ONE large pre-drawn
% display the student reads or names, then a bold prompt and a two-across grid of numbered
% problems with 1.3–1.8cm of writing room. Problems run 1..14 (rows 1–4: 1–2, 3–4, 5–6, 7–10).
% DENSITY: no \blank{} at all this lesson — the skip verdict for 951 is a \labelbox name under
% the row-2 display; no mid-sentence blanks; every display is pre-drawn; the page belongs to
% the student's pen. (\wordbank is defined but unused; a fill table was cut 2026-09-14, and on
% 2026-09-14 the displays were tightened and rows 3's four problems merged into two, to hold
% the notes to four pages.)
%   I DO   : the teacher reads the definition, marks up the display, works the first problem
%            of each grid (1, 3, 5, 7) and reads the digit strip aloud, stopping at 951.
%   WE DO  : the class works the rest of each grid, students holding the pen; the last row,
%            Guided Practice (11–14), is worked entirely together on Cedar Ridge Apartments
%            instead of Riverbend.
%   YOU DO : nothing on this page — ../homework/ IS the individual practice, started in the
%            last ten minutes of the period. No independent set, no reflection box, no
%            objectives box (the cover carries the targets).
% THE TRAP is problem 2 (the teacher's own third period — "random" does not mean "unplanned";
% the test is WHO COULD NEVER HAVE BEEN PICKED).
% THE CRUX is problem 9 (row 4, "Cluster"): Maya's two random homerooms are 60 ninth graders
% and zero seniors, nobody counted wrong — the hook vote is re-taken there (PS.DC.2a–b).
% Problem 13 is the crux transferred to Cedar Ridge; problem 14 is the choose-and-justify
% (PS.DC.2b) plus the list-of-buildings question.
%
% THE DATA
%   Riverbend HS: 900 students, the council surveys 60. Grades 270/240/210/180 -> at 1/15:
%     18/16/14/12. Digit strip 7 1 4 0 2 8 9 5 1 3 6 6 -> students 714, 28, (skip 951), 366.
%     Systematic k = 900/60 = 15, start 7 -> 7, 22, 37, 52 (the warm-up's numbers).
%     30 homerooms of 30: homerooms 1-9 hold the 9th grade, 10-17 the 10th, 18-24 the 11th,
%     25-30 the 12th. Maya draws homerooms 4 and 7 -> 2 x 30 = 60 students, all 9th graders.
%     About one draw in four (100/435 = 23%) lands both rooms in a single grade (C(30,2)=435;
%     same-grade pairs 36+28+21+15 = 100).
%   Guided Practice, Cedar Ridge Apartments: 600 households in 30 buildings of 20; buildings
%     1-6 studios (120), 7-20 one-bedrooms (280), 21-30 two-bedrooms (200); 60 households at
%     1/10 -> 12/28/20; a cluster draw of buildings 2, 5, 6 -> 60 studio households.
%
% PAGE LOCKSTEP with notes_key/: the key is GENERATED from this file by templates/lesson/mkkey.py
% (spec: templates/lesson/examples/lesson03_notes_key.json) — every \blank{} becomes an \ans{},
% every \vterm{} a \vtermans{}{}, every \par\writelines{n} n \ansline{} calls, and the answer
% argument of every \pcell, \writespace and \labelbox is filled. Answer spaces are fixed
% heights, so the two files cannot drift. KEEP EVERY \pcell ON ONE SOURCE LINE — mkkey matches
% line by line.
% TABLE MECHANICS: inside a cell, `\\` ENDS THE ROW — break lines with \par. A row cannot break
% across pages: the figure gets its own sub-row (`\\` then `& ...`) and EACH GRID ROW is its own
% sub-row — close the probgrid, `\\ &`, reopen as probgrid* (no top rule).
\documentclass[12pt]{article}
\usepackage{saar-article}
\usepackage{saar-boxes}
\newcommand{\TallMath}[1]{$\displaystyle #1\rule[-1.4em]{0pt}{3.2em}$}
% Word bank strip — ONLY above a table the student fills. Defined per document (copied from
% 1.2) and currently unused here; the key inherits it either way.
\newcommand{\wordbank}[1]{\par\vspace{2pt}\noindent\fcolorbox{goldacc}{goldbg}{\parbox{\dimexpr\linewidth-2\fboxsep-2\fboxrule\relax}{\small\textbf{Word bank:} #1}}\par\vspace{4pt}}
% Vocabulary rows: a FIXED-HEIGHT row carrying the term label and OPEN writing space —
% no underline beside the term and no rule line (user correction 2026-09-06: the black
% answer line crowded the row; students write beside and below the term). The key fills
% exactly the same height, so the box cannot drift blank vs. keyed. saar-article's
% \termblank adds an inline blank and a rule line, so the pair is defined here (the key is
% generated from this file and inherits both). 2.3cm per row gives students three
% handwritten lines of room; keep key definitions to two lines.
\newcommand{\termrowheight}{2.3cm}
\newlength{\termrowinset}\setlength{\termrowinset}{7pt}
\newcommand{\vterm}[1]{%
  \par\noindent\begin{minipage}[t][\termrowheight][t]{\linewidth}%
    \vspace*{\termrowinset}%
    \textbf{\textcolor{cerulean}{#1:}}%
  \end{minipage}\par}
\newcommand{\vtermans}[2]{%
  \par\noindent\begin{minipage}[t][\termrowheight][t]{\linewidth}%
    \vspace*{\termrowinset}%
    \textbf{\textcolor{cerulean}{#1:}}\quad\ans{#2}%
  \end{minipage}\par}

\begin{document}

\pageheader{Unit 1, Lesson 1.3}{Guided Notes \& Practice}
% NAMESTRIP: no \namedateperiod here. The student writes their name once, on the
% cover this component is stapled behind.

\vspace{0.15in}

\begin{vocabbox}
\small
Fill in each term \textbf{as we name it} in the notes below.
\vterm{Sampling frame}
\vterm{Simple random sample (SRS)}
\vterm{Stratified sample}
\vterm{Systematic sample}
\vterm{Cluster sample}
\end{vocabbox}

\boxguard[12]
\begin{hookbox}
\small
The council needs $60$ of Riverbend's $900$ students. Maya puts the $30$ homeroom numbers in a
hat, draws two --- rooms $4$ and $7$ --- and surveys all $30$ students in each. Pure chance.

\vspace{3pt}
\textbf{Is Maya's sample as trustworthy as drawing $60$ student names out of the hat?} Circle
one: \quad yes \qquad no \qquad it depends \quad --- and say why you think so. We vote again
at problem~9.
\par\writelines{2}
\end{hookbox}

\small
\begin{guidednotes}
% ── ROW 1: the frame, and who chooses ────────────────────────────────────────
\mainidea[Chance or convenience]{Who Chooses?} &
The \textbf{sampling frame} is the list a study can choose from. In a \textbf{probability
sample} chance chooses; in a \textbf{convenience sample} the researcher takes whoever is
easiest to reach. \textbf{The test: who could never have been picked?}
\\
& {\centering
\begin{tikzpicture}[scale=0.68]
  % the frame: the office list (node fonts do NOT scale with `scale=`, so every
  % label here is \tiny and short enough to sit inside its own box)
  \fill[frost, rounded corners=4pt] (0,2.85) rectangle (11.5,3.75);
  \draw[cerulean, thick, rounded corners=4pt] (0,2.85) rectangle (11.5,3.75);
  \node[font=\tiny\bfseries\color{cerulean}] at (5.75,3.52) {THE OFFICE LIST --- all $900$ students, numbered $001$--$900$};
  \foreach \i in {0,...,21} { \fill[slate] ({1.7+\i*0.38},3.12) circle (0.06); }
  % left panel: convenience (the pool's Tuesday evening, from 1.2)
  \fill[goldbg, rounded corners=4pt] (0,0) rectangle (5.4,2.3);
  \draw[goldacc, thick, rounded corners=4pt] (0,0) rectangle (5.4,2.3);
  \node[font=\tiny\bfseries\color{goldacc}] at (2.7,2.05) {THE RESEARCHER CHOOSES};
  \node[font=\tiny, align=center] at (2.7,1.32) {``Ask whoever shows up\\ Tuesday evening.''\\ the schedule chose};
  \node[font=\tiny\itshape, align=center] at (2.7,0.42) {who could never be picked?\\ anyone who swims another night};
  % right panel: chance
  \fill[frost, rounded corners=4pt] (6.1,0) rectangle (11.5,2.3);
  \draw[cerulean, thick, rounded corners=4pt] (6.1,0) rectangle (11.5,2.3);
  \node[font=\tiny\bfseries\color{cerulean}] at (8.8,2.05) {CHANCE CHOOSES};
  \node[font=\tiny, align=center] at (8.8,1.32) {number the list, let a\\ generator draw $60$\\ every name had a chance};
  \node[font=\tiny\itshape, align=center] at (8.8,0.42) {who could never be picked?\\ nobody on the list};
  \draw[->, very thick, goldacc] (2.7,2.38) -- (2.7,2.8);
  \draw[->, very thick, cerulean] (8.8,2.38) -- (8.8,2.8);
\end{tikzpicture}\par\vspace{4pt}
Frame: \labelbox{3.0cm}{}\quad Left: \labelbox{2.1cm}{}\quad Right: \labelbox{2.1cm}{}\par}
\\
& \notesprompt{Who chose each sample --- chance, or convenience? Name someone who could never have been picked.}
\begin{probgrid}{|Y|Y|}
\pcell{1}{A pool staff member asks the $75$ members who arrive on Tuesday evening. Then: all $900$ students are numbered and a computer picks $60$. Who chose each?}{1.6cm}{} & \pcell{2}{A teacher surveys her own third-period class --- she picked no favorites. \textbf{Vote first, then decide.}}{1.6cm}{} \\ \hline
\end{probgrid}
\\ \hline
% ── ROW 2: the two list techniques ───────────────────────────────────────────
\mainidea[Number the frame, then draw]{SRS \& Systematic} &
In a \textbf{simple random sample} of size $n$, every individual --- and every group of $n$
--- has an equal chance. \stepnum{1}\ Number the frame $1$ to $N$. \stepnum{2}\ Read the
generator's digits in groups as long as $N$. \stepnum{3}\ \textbf{Skip} any group out of range
or already used.\par
From an ordered list, a \textbf{systematic sample} takes every $k$th individual after a random
\textbf{start}, where $k = N \div n$.
\\
& {\centering
\begin{tikzpicture}[scale=0.60]
  % top: the digit strip, read three at a time
  \node[anchor=west, font=\tiny\bfseries\color{cerulean}] at (0,3.35) {THE GENERATOR'S DIGITS, read three at a time --- $900$ has three digits};
  \foreach \g/\a/\b/\c/\lab in {0/7/1/4/{student 714}, 1/0/2/8/{student 28}, 2/9/5/1/{}, 3/3/6/6/{}} {
    \draw[fill=white, thick] ({\g*2.9},2.45) rectangle ++(0.7,0.7);
    \draw[fill=white, thick] ({\g*2.9+0.7},2.45) rectangle ++(0.7,0.7);
    \draw[fill=white, thick] ({\g*2.9+1.4},2.45) rectangle ++(0.7,0.7);
    \node[font=\ttfamily\large] at ({\g*2.9+0.35},2.8) {\a};
    \node[font=\ttfamily\large] at ({\g*2.9+1.05},2.8) {\b};
    \node[font=\ttfamily\large] at ({\g*2.9+1.75},2.8) {\c};
    \node[font=\tiny\bfseries\color{goldacc}] at ({\g*2.9+1.05},2.17) {\lab};
  }
  % bottom: the ordered list with every 15th name
  \node[anchor=west, font=\tiny\bfseries\color{cerulean}] at (0,1.80) {THE LIST IN ORDER --- a random start, then every $k$th name};
  \draw[thick] (0.2,0.80) -- (11.2,0.80);
  \foreach \i in {1,...,60} { \draw ({0.2+\i*0.18},0.72) -- ({0.2+\i*0.18},0.88); }
  \foreach \p in {7,22,37,52} { \fill[goldacc] ({0.2+\p*0.18},0.80) circle (0.12); \node[below, font=\tiny\bfseries] at ({0.2+\p*0.18},0.63) {\p}; }
  \node[font=\tiny, anchor=west] at (11.25,0.80) {\dots\ $900$};
  \draw[<->, goldacc, thick] ({0.2+7*0.18},1.22) -- ({0.2+22*0.18},1.22);
  \node[right, font=\tiny\bfseries\color{goldacc}] at ({0.2+22*0.18+0.18},1.22) {$k = 900 \div 60 = 15$};
\end{tikzpicture}\par\vspace{4pt}
Top: \labelbox{2.6cm}{}\quad Bottom: \labelbox{2.2cm}{}\quad $951$: \labelbox{2.0cm}{}\par}
\\
& \notesprompt{Run both techniques on Riverbend.}
\begin{probgrid}{|Y|Y|}
\pcell{3}{Students are numbered $001$--$900$. What do you do with $951$ --- \emph{skip} it, or \emph{stop}? What does the generator do next?}{1.4cm}{} & \pcell{4}{Systematic: $k = 900 \div 60 = {}$? The random start lands on $7$. Write the next three names taken --- and circle the one number chance chose.}{1.4cm}{} \\ \hline
\end{probgrid}
\\ \hline
% ── ROW 3: stratified — some from every group ────────────────────────────────
\mainidea[Some from every group]{Stratified} &
\textbf{Strata} are groups \textbf{alike inside} and \textbf{different from group to group}
(one is a \emph{stratum}). A \textbf{stratified sample} takes a separate random sample from
\emph{every} stratum, sized to its share --- so it \textbf{guarantees} every grade is heard,
in proportion. Seniors drive, so grade is a natural set of strata.
\\
& {\centering
\begin{tikzpicture}[scale=0.66]
  % the school as four strata, widths in proportion to 270/240/210/180
  \node[anchor=west, font=\tiny\bfseries\color{cerulean}] at (0,3.25) {ALL $900$ --- four strata, then $60/900 = 1/15$ of \emph{every} one};
  \foreach \x/\w/\g/\c/\col in {0/3.3/9th/270/frost, 3.45/2.93/10th/240/goldbg, 6.53/2.57/11th/210/greenbg, 9.25/2.2/12th/180/plumbg} {
    \fill[\col, rounded corners=3pt] (\x,1.95) rectangle ++(\w,1.0);
    \draw[cerulean, thick, rounded corners=3pt] (\x,1.95) rectangle ++(\w,1.0);
    \node[font=\tiny\bfseries] at ({\x+\w/2},2.68) {\g};
    \node[font=\tiny] at ({\x+\w/2},2.25) {$\c$ students};
    \draw[->, thick, charcoal] ({\x+\w/2},1.88) -- ({\x+\w/2},1.25);
    \node[font=\tiny, right] at ({\x+\w/2+0.05},1.58) {$\div 15$};
    \draw[goldacc, thick, fill=white, rounded corners=3pt] ({\x+\w/2-0.5},0.42) rectangle ++(1.0,0.78);
  }
  \node[font=\small\bfseries] at (1.65,0.81) {$18$};
  \node[font=\small\bfseries] at (4.915,0.81) {?};
  \node[font=\small\bfseries] at (7.815,0.81) {?};
  \node[font=\small\bfseries] at (10.35,0.81) {?};
  \node[font=\tiny\itshape, anchor=west] at (0,0.05) {the sample: some from every grade, in proportion --- $60$ in all};
\end{tikzpicture}\par\vspace{4pt}
Each grade is one \labelbox{2cm}{}\quad it reaches \labelbox{1.5cm}{} grades\par}
\\
& \notesprompt{Finish the allocation the warm-up started.}
\begin{probgrid}{|Y|Y|}
\pcell{5}{The warm-up found $270 \div 15 = 18$ ninth graders. Compute the other three shares --- then add all four. Do they total $60$?}{1.6cm}{} & \pcell{6}{Stratifying \emph{guarantees} every grade is heard; a fair SRS of $60$ names only makes it \emph{likely}. Why?}{1.6cm}{} \\ \hline
\end{probgrid}
\\ \hline
% ── ROW 4: THE CRUX — cluster, and Maya's homeroom draw ──────────────────────
\mainidea[Everyone from a few groups]{Cluster} &
When there is no list of individuals, only a list of groups, a \textbf{cluster sample} draws a
few whole groups at random and collects data from \emph{everyone} in them.
\textbf{It works only when each cluster is a small mix of the whole.} Riverbend's $900$
students sit in $30$ homerooms of $30$; Maya drew rooms $4$ and $7$.
\\
& {\centering
\begin{tikzpicture}[scale=0.62]
  \node[anchor=west, font=\tiny\bfseries\color{cerulean}] at (0,3.62) {RIVERBEND'S $30$ HOMEROOMS OF $30$ --- Maya's draw is circled};
  \foreach \n in {1,...,9}   { \fill[frost]   ({mod(\n-1,10)*1.1},{2.4-int((\n-1)/10)*1.0}) rectangle ++(1.0,0.75); }
  \foreach \n in {10,...,17} { \fill[goldbg]  ({mod(\n-1,10)*1.1},{2.4-int((\n-1)/10)*1.0}) rectangle ++(1.0,0.75); }
  \foreach \n in {18,...,24} { \fill[greenbg] ({mod(\n-1,10)*1.1},{2.4-int((\n-1)/10)*1.0}) rectangle ++(1.0,0.75); }
  \foreach \n in {25,...,30} { \fill[plumbg]  ({mod(\n-1,10)*1.1},{2.4-int((\n-1)/10)*1.0}) rectangle ++(1.0,0.75); }
  \foreach \n in {1,...,30}  {
    \draw[cerulean] ({mod(\n-1,10)*1.1},{2.4-int((\n-1)/10)*1.0}) rectangle ++(1.0,0.75);
    \node[font=\scriptsize\bfseries] at ({mod(\n-1,10)*1.1+0.5},{2.4-int((\n-1)/10)*1.0+0.375}) {\n}; }
  \draw[redacc, very thick] (3.8,2.775) circle (0.55);
  \draw[redacc, very thick] (7.1,2.775) circle (0.55);
  % legend --- one entry every 2.9 units; keep each label under ~14 characters
  \foreach \x/\col/\lab in {0/frost/{9th: 1--9}, 2.9/goldbg/{10th: 10--17}, 5.8/greenbg/{11th: 18--24}, 8.7/plumbg/{12th: 25--30}} {
    \fill[\col] (\x,-0.2) rectangle ++(0.32,0.28); \draw[cerulean] (\x,-0.2) rectangle ++(0.32,0.28);
    \node[font=\tiny, anchor=west] at ({\x+0.38},-0.06) {\lab}; }
\end{tikzpicture}\par}
\\
& \textbf{``Chance chose'' is not enough --- the technique has to match the groups.} Alike
inside $\Rightarrow$ some from \emph{every} group (stratified); a small mix of the whole
$\Rightarrow$ everyone from a \emph{few} (cluster).\par
\textbf{In your own words:} why can a random draw of two homerooms miss every senior?
\writespace{1.0cm}{}
\\
& \notesprompt{Now settle the hook vote.}
\begin{probgrid}{|Y|Y|}
\pcell{7}{Maya drew rooms $4$ and $7$ and surveyed everyone in both. How many students is that, and which technique is it?}{1.3cm}{} & \pcell{8}{Now read the map. Which grade is room $4$? Room $7$? How many seniors are in Maya's $60$?}{1.3cm}{} \\ \hline
\end{probgrid}
\\
& \begin{probgrid}*{|Y|Y|}
\pcell{9}{\textbf{Vote again:} yes \; no \; it depends. Did chance choose? Did anyone count wrong? Then what went wrong?}{1.7cm}{} & \pcell{10}{Is a Riverbend homeroom a stratum or a cluster? Which technique matches groups like that, and what does it take from each?}{1.7cm}{} \\ \hline
\end{probgrid}
\\ \hline
% ── ROW 5: GUIDED PRACTICE — we do, together, a fresh context ────────────────
\mainidea[Guided practice]{Cedar Ridge} &
Cedar Ridge Apartments has $600$ households in $30$ buildings of $20$, and management wants to
survey $60$ about parking. Buildings $1$--$6$ are studios ($120$), $7$--$20$ one-bedrooms
($280$), $21$--$30$ two-bedrooms ($200$) --- and two-bedroom households own the most cars.
\\
& {\centering
\begin{tikzpicture}[scale=0.64]
  \node[anchor=west, font=\tiny\bfseries\color{cerulean}] at (0,3.62) {CEDAR RIDGE --- $30$ buildings of $20$; the hat drew $2$, $5$, $6$};
  \foreach \n in {1,...,6}   { \fill[frost]   ({mod(\n-1,10)*1.1},{2.4-int((\n-1)/10)*1.0}) rectangle ++(1.0,0.75); }
  \foreach \n in {7,...,20}  { \fill[goldbg]  ({mod(\n-1,10)*1.1},{2.4-int((\n-1)/10)*1.0}) rectangle ++(1.0,0.75); }
  \foreach \n in {21,...,30} { \fill[greenbg] ({mod(\n-1,10)*1.1},{2.4-int((\n-1)/10)*1.0}) rectangle ++(1.0,0.75); }
  \foreach \n in {1,...,30}  {
    \draw[cerulean] ({mod(\n-1,10)*1.1},{2.4-int((\n-1)/10)*1.0}) rectangle ++(1.0,0.75);
    \node[font=\scriptsize\bfseries] at ({mod(\n-1,10)*1.1+0.5},{2.4-int((\n-1)/10)*1.0+0.375}) {\n}; }
  \draw[redacc, very thick] (1.6,2.775) circle (0.55);
  \draw[redacc, very thick] (4.9,2.775) circle (0.55);
  \draw[redacc, very thick] (6.0,2.775) circle (0.55);
  % legend --- the counts live in the printed sentence above, so the labels stay short
  \foreach \x/\col/\lab in {0/frost/{studios: 1--6}, 3.7/goldbg/{one-bed: 7--20}, 7.4/greenbg/{two-bed: 21--30}} {
    \fill[\col] (\x,-0.2) rectangle ++(0.32,0.28); \draw[cerulean] (\x,-0.2) rectangle ++(0.32,0.28);
    \node[font=\tiny, anchor=west] at ({\x+0.38},-0.06) {\lab}; }
\end{tikzpicture}\par}
\\
& \notesprompt{We work these together.}
\begin{probgrid}{|Y|Y|}
\pcell{11}{Name the technique in each plan: (i)~a computer draws $60$ of the numbered $600$; (ii)~$3$ buildings from a hat, all $20$ in each; (iii)~a random sample from each unit type, in proportion; (iv)~every $10$th name on the resident list after a random start.}{1.8cm}{} & \pcell{12}{The stratified plan: $60$ of $600$ is $1/10$. Compute all three shares and check the total.}{1.8cm}{} \\ \hline
\end{probgrid}
\\
& \begin{probgrid}*{|Y|Y|}
\pcell{13}{The hat gives buildings $2$, $5$, and $6$. Who is in the sample? Is a building a mix of the complex, or a stratum? What did the draw get wrong, even though chance chose?}{1.7cm}{} & \pcell{14}{Which plan should management run --- \emph{because the groups are}\,\dots? And with only a list of \emph{buildings}, never of residents, which technique still runs?}{1.7cm}{} \\ \hline
\end{probgrid}
\\ \hline
\end{guidednotes}

% ── THE NOTES END AT GUIDED PRACTICE ─────────────────────────────────────────
% There is deliberately no "you do" here: ../homework/ is the individual practice,
% started in class in the last ten minutes while the teacher circulates.

\end{document}
